% Standalone mathematical note, accompanying Appendix A.2.
\paragraph{Two-context no-reversal result.}
Assume a finite support with at most two contexts, probabilities $p_i>0$,
conditional exposures $w_i>0$, and finite nonnegative conditional losses
$\ell_i$. Acceptance probabilities $a_i\in[0,1]$ are unrestricted.
The only constraints are $\sum_i p_i a_i(\ell_i-rw_i)\leq0$,
$\sum_i p_i a_i\geq c_0$, and positive accepted exposure.
Whenever this feasible set is nonempty, the case and exposure objectives
have the same unique optimum in acceptance probabilities.

Write $g_i=p_i(\ell_i-rw_i)$. Both objectives are strictly increasing
in every coordinate because their coefficients are respectively $p_i>0$
and $p_iw_i/T>0$, where $T=\sum_i p_iw_i$.
For one context, feasibility implies $g_1\leq0$, so $a_1=1$ is optimal.
For two contexts, if both $g_i\leq0$, $(1,1)$ is feasible whenever any gate
is feasible and dominates every gate. If both are positive, no gate with
positive accepted exposure is feasible. Otherwise relabel so that
$g_1\leq0<g_2$. The componentwise greatest feasible candidate is
\[
 a_1^*=1,\qquad a_2^*=\min\{1,-g_1/g_2\}.
\]
Indeed, any feasible gate satisfies
$a_2\leq(-g_1/g_2)a_1\leq-g_1/g_2$ and $a_i\leq1$.
The candidate meets the pooled risk constraint; since it dominates every
feasible gate, it also meets the row floor whenever any feasible gate does.
Its positive accepted exposure follows from $p_1w_1>0$.
Strict coordinate monotonicity establishes uniqueness for both objectives.

Thus strict case/exposure reversal under these constraints requires at least
three positive-probability, positive-exposure contexts. Uniqueness refers to
the acceptance probabilities, not to the random seeds realizing them.
This result is not asserted for additional constraints or a restricted
acceptance-policy family. The existing finite-grid verifier is a numerical
check of examples; the argument above proves the continuous-moment statement.
